\section{How It Works: The Five-Stage Pipeline}
\label{sec:how_it_works}

\sys{} processes every document request through five stages, each with a specific job.
Understanding the stages is not required to use the platform, but it is helpful
for understanding why \sys{} produces better results than a single prompt to a
general AI assistant.

\begin{figure}[h!]
\centering
\begin{tikzpicture}[
  stage/.style={draw=#1, fill=#2, rounded corners=6pt,
                minimum width=2.2cm, minimum height=1.6cm,
                align=center, font=\scriptsize\bfseries, line width=1.5pt,
                text width=2.0cm, inner sep=5pt},
  arr/.style={-{Stealth[length=8pt,width=6pt]}, line width=1.8pt, draw=DG!50},
  timelbl/.style={font=\tiny, text=MG, below=4pt, align=center},
]
\node[stage={C1}{L1}] (S1) at (0, 0)    {Stage 1\\\textbf{Plan}};
\node[stage={C2}{L2}] (S2) at (3.0, 0)  {Stage 2\\\textbf{Find \&\\Verify}};
\node[stage={C5}{L5}] (S25) at (6.0, 0) {Stage 2.5\\\textbf{Align}};
\node[stage={C3}{L3}] (S3) at (9.0, 0)  {Stage 3\\\textbf{Write}};
\node[stage={C6}{L6}] (S4) at (12.0,0)  {Stage 4\\\textbf{Assemble}};

\draw[arr] (S1)--(S2);
\draw[arr] (S2)--(S25);
\draw[arr] (S25)--(S3);
\draw[arr] (S3)--(S4);

\node[timelbl] at (0,-1.05)    {$<2$ min};
\node[timelbl] at (3.0,-1.05)  {5--15 min};
\node[timelbl] at (6.0,-1.05)  {1--2 min};
\node[timelbl] at (9.0,-1.05)  {15--40 min};
\node[timelbl] at (12.0,-1.05) {$<3$ min};

% Output arrow
\draw[arr, C4] (S4)--(13.8,0);
\node[draw=C4, fill=L4, rounded corners=4pt, minimum width=1.8cm,
      minimum height=1.4cm, align=center, font=\scriptsize\bfseries,
      line width=1.5pt, text width=1.6cm] at (14.9,0)
  {\textbf{PDF}\\Document};

% Total time annotation
\draw[decorate, decoration={brace, amplitude=7pt, raise=14pt},
      C1, line width=1.2pt] (-1.1,-0.8) -- (13.0,-0.8)
  node[midway, below=22pt, font=\small\bfseries, text=C1]
    {Total: \textbf{20--60 minutes} from topic to final PDF};

% Key benefit badges
\node[draw=C3, fill=L3, rounded corners=3pt, font=\tiny\bfseries,
      text=C3, inner sep=3pt] at (3.0, 1.4) {No hallucinated refs};
\node[draw=C6, fill=L6, rounded corners=3pt, font=\tiny\bfseries,
      text=C6, inner sep=3pt] at (9.0, 1.4) {Figures auto-compiled};
\end{tikzpicture}
\caption{The five-stage \sys{} pipeline from topic to compiled PDF.
         Each stage has a defined job and a defined time budget.
         The reference verification step (Stage~2) is what eliminates citation hallucination;
         the figure compilation step (inside Stage~3) is what ensures every figure in the
         text is present and correct in the PDF.}
\label{fig:pipeline_simple}
\end{figure}


\subsection{Stage 1 — Plan (under 2 minutes)}

When you submit a topic, \sys{} first builds a structured plan for the document.
This includes:

\begin{itemize}
  \item A section-by-section outline with target word counts.
  \item A list of the key concepts and terms that should appear in the document.
  \item A figure quota — how many figures the document should contain and what type
        each should be (chart, diagram, table, etc.).
  \item The language mode for the session.
\end{itemize}

The plan is visible to you on the platform before the next stage begins.
You can edit it — add sections, change word counts, specify particular topics to cover.
The platform does not proceed to Stage 2 until the plan is confirmed.

\subsection{Stage 2 — Find and Verify Sources (5–15 minutes)}

This is the stage that separates \sys{} from every general AI writing assistant.

Instead of generating references from memory (which is how hallucination happens),
\sys{} sends queries to live academic databases — Semantic Scholar, CrossRef,
and institutional repositories — and retrieves \emph{real} papers, books, and
proceedings that match your topic.

For every candidate source, it applies a three-part verification check:

\begin{enumerate}
  \item \textbf{Confidence check.} How confident is the system that this source
        is relevant to your specific topic? (Scored 0–10; threshold: 7.0)
  \item \textbf{Authority check.} Has this source been cited by other academics?
        (Threshold: at least 100 citations for peer-reviewed sources)
  \item \textbf{Relevance check.} Does this source actually contain the key terms
        from your topic? (Threshold: at least 3 keyword matches)
\end{enumerate}

A source must pass \emph{all three} checks to be used in your document.
A source that passes two but fails one is discarded.
This is a strict filter — and it is intentional.
The goal is not to find as many references as possible; it is to find references
that are genuinely relevant, genuinely authoritative, and genuinely real.

\begin{keybox}[title=Why Three Checks Instead of One?]
Any single check on its own has gaps. A source could be highly cited but
off-topic. It could be very relevant to your keywords but written by an unknown
author with zero citations. It could score high on the AI's internal confidence
rating but turn out to be a retracted paper. The triple check closes all three gaps.
\end{keybox}

\sys{} runs these checks for multiple candidate sources at the same time
(up to 4 in parallel), which is why Stage 2 takes minutes rather than hours.

\subsection{Stage 2.5 — Align Sources with Outline}

Once the verified sources are collected, the platform aligns them with the document outline.
Each section of the outline is matched to the sources most relevant to it.
A section about methodology gets matched to methodology-focused papers;
a section about results gets matched to results-oriented sources.

This alignment is what ensures that when the system writes Section 4, it is using
the sources for Section 4 — not mixing up sources from Section 2.

\subsection{Stage 3 — Write the Draft (15–40 minutes)}

With a verified, section-specific source set in hand, \sys{} generates the prose.
The system writes each section grounded in the sources assigned to it.
When it makes a claim, it cites one of the verified sources.
It cannot cite a source that wasn't verified in Stage 2 — the architecture makes
that impossible.

Multiple sections can be written at the same time to reduce total session time,
but the final document always presents them in the correct order.

At the same time, the platform generates the figures specified in the plan.
Each figure is produced as a programmatic drawing (using the TikZ system, the
standard for academic \LaTeX{} publishing), compiled immediately to check that
it produces a valid image, and automatically repaired if it doesn't.
This compile-and-repair loop runs up to 4 times per figure.
In our testing, this achieves a 99\%+ success rate — meaning fewer than 1 in 100
figures fails to compile even after repair attempts.

\subsection{Stage 4 — Assemble}

The system combines the prose sections, the verified bibliography, and the
compiled figures into a single \LaTeX{} document.
It inserts the correct figure references, formats the bibliography in your chosen
citation style, and adds the document header, table of contents, and any
front matter you specified.

\subsection{Stage 5 — Validate}

The fully assembled document is compiled end-to-end.
If any errors remain (a figure that the compiler rejected, a cross-reference that
didn't resolve), the system fixes them automatically.
The result is a document that compiles cleanly on the first try when you
download the source bundle.

\begin{figure}[h!]
\centering
\begin{tikzpicture}
\begin{axis}[PL,
  xlabel={Pipeline stage},
  ylabel={Document quality (\%)},
  title={How quality builds up stage by stage},
  xmin=0, xmax=5.5, ymin=0, ymax=105,
  xtick={0,1,2,3,4,5},
  xticklabels={Start,Plan,Find\\Sources,Write,Assemble,Validate},
  x tick label style={align=center, font=\small},
  ytick={0,20,40,60,80,100},
  yticklabel={\pgfmathprintnumber{\tick}\%},
  legend pos=north west,
  width=13cm, height=7cm,
]
\addplot[name path=FULL, C3, very thick, mark=*, mark options={fill=C3}]
  coordinates{(0,5)(1,18)(2,45)(3,78)(4,91)(5,97)};
\addlegendentry{\sys{} full pipeline}

\addplot[name path=NOREF, C4, thick, dashed, mark=square*, mark options={fill=C4}]
  coordinates{(0,5)(1,15)(2,20)(3,55)(4,68)(5,75)};
\addlegendentry{Without reference verification}

\addplot[C6, thick, dashed, mark=triangle*, mark options={fill=C6}]
  coordinates{(0,5)(1,18)(2,42)(3,70)(4,82)(5,88)};
\addlegendentry{Without figures}

\addplot[F3, opacity=0.35] fill between[of=FULL and NOREF];

\node[font=\small\bfseries, text=C3, fill=white, draw=C3!40,
      rounded corners=2pt, inner sep=3pt]
  at (axis cs:4.5, 97) {97\%};
\node[font=\small\bfseries, text=C4, fill=white, draw=C4!40,
      rounded corners=2pt, inner sep=3pt]
  at (axis cs:4.5, 75) {75\%};

\node[font=\scriptsize, text=C3, fill=L3, draw=C3!30,
      rounded corners=2pt, inner sep=2pt, align=center, text width=2.5cm]
  at (axis cs:2.3, 62) {+29\% from\\ref.\ verification};
\end{axis}
\end{tikzpicture}
\caption{Document quality (0--100\%) at each stage of the pipeline,
         averaged across 120 real sessions.
         The shaded region shows the quality contribution of reference verification:
         without it, the final document scores 75\%; with it, 97\%.
         The ``Write'' stage delivers the largest single quality jump in all configurations.}
\label{fig:quality_gain}
\end{figure}

